\section{Related Work}
\label{sec:related}

\subsection{Bayesian dynamic borrowing and the exchangeability assumption}

Integrating external information into a concurrent study has a long history
in clinical biostatistics, beginning with the hierarchical framework of
\citet{pocock1976combination} for combining randomized and historical
controls.
Let $D_c$ denote the data from a small current study and $D_e$ the data from
an external or historical source, with a current-study target parameter
$\theta$.
The central statistical assumption underlying any borrowing procedure is
\emph{exchangeability} between the two sources: that, conditional on
covariates, observations from $D_c$ and $D_e$ share a common generative
mechanism.
When exchangeability holds, pooling reduces variance and improves
precision, yielding effective sample-size gains that are especially valuable
in rare-disease, pediatric, and oncology trials where concurrent accrual is
limited \citep{berry2013bayesian,ventz2019design,edwards2024using}.
When it fails---because of secular drift, differing diagnostic criteria, or
changes in the outcome mechanism---naive pooling induces bias and can inflate
type~I error \citep{viele2014use,koppschneider2020power,psioda2019bayesian}.
This tension motivates \emph{dynamic} rather than fixed borrowing: methods
that quantify cross-source compatibility and modulate the contribution of
$D_e$ accordingly.
Existing dynamic borrowing methods mainly differ in how they assess
compatibility between $D_c$ and $D_e$ and how that assessment is translated
into borrowing strength.

\paragraph{Power prior.}
The power prior \citep{ibrahim2000power,ibrahim2015power} incorporates
external information by discounting the external likelihood through a scalar
weight $a_0 \in [0,1]$:
\begin{equation}
p(\theta \mid D_c, D_e, a_0)
\;\propto\;
L(\theta; D_c)\,L(\theta; D_e)^{a_0}\,\pi(\theta).
\end{equation}
When $a_0 = 1$ the two sources are fully pooled, whereas $a_0 = 0$ ignores
the external data entirely.
Fixing $a_0$ a priori is difficult in practice.
The normalized power prior \citep{duan2006evaluating,ye2022normalized} treats
$a_0$ as unknown and integrates over it, and empirical-Bayes adaptive
variants \citep{gravestock2017adaptive} plug in a data-driven estimate.

\paragraph{Commensurate prior.}
The commensurate prior \citep{hobbs2011commensurate,hobbs2012commensurate}
introduces separate parameters for the current and external sources and
links them through a precision parameter $\tau$:
\begin{align}
\theta_e &\sim \pi(\theta_e), &
\theta_c \mid \theta_e, \tau &\sim \mathcal{N}(\theta_e,\tau^{-1}).
\end{align}
Large $\tau$ encourages $\theta_c$ and $\theta_e$ to be close, whereas small
$\tau$ allows them to differ.
With an appropriate hyperprior on $\tau$, the posterior adaptively reduces
borrowing under prior--data conflict \citep{evans2006checking,
presanis2013conflict}.

\paragraph{Meta-analytic-predictive and robust MAP priors.}
The meta-analytic-predictive (MAP) prior
\citep{neuenschwander2010summarizing} derives an informative prior for
$\theta_c$ by treating historical trials as exchangeable draws from a common
hierarchical model.
Because exchangeability can fail in ways that are hard to anticipate,
\citet{schmidli2014robust} proposed a robust MAP prior that mixes the
informative component with a vague one:
\begin{equation}
\pi_{\text{robust}}(\theta_c)
=
(1-\epsilon)\,\pi_{\text{MAP}}(\theta_c)
+
\epsilon\,\pi_{\text{vague}}(\theta_c),
\end{equation}
where $\epsilon$ controls protection against conflict between historical and
current data.
This construction is attractive in practice because $\epsilon$ has a
transparent interpretation, but the choice of $\epsilon$ and of the vague
component materially affects the resulting borrowing behavior.

\paragraph{Limitations of classical BDB.}
Several studies have documented recurring practical limitations of these
methods.
First, operating characteristics are sensitive to the discount
hyperparameters $(a_0,\tau,\epsilon)$, and strict type~I error control
essentially eliminates the power gains that motivated borrowing in the first
place \citep{koppschneider2020power,galwey2017supplementation,
vanrosmalen2018including}.
Second, classical BDB methods summarize cross-source compatibility through
\emph{one or a few low-dimensional quantities}, which can mask clinically
important source differences: in particular, covariate shift and outcome
drift enter the same scalar $a_0$ or $\tau$, even though they have very
different implications for bias in the current-study estimand
\citep{viele2014use,dahabreh2019generalizing}.
Third, these models require repeated posterior fitting---typically via
MCMC---for each new dataset pair, which is expensive when operating
characteristics must be simulated across many scenarios, when borrowing
decisions must be reassessed for several candidate external cohorts, or when
the analysis must be rerun as data accrue.
These limitations motivate borrowing procedures that can (i) exploit richer,
multi-dimensional summaries of source compatibility, and (ii) amortize
inference cost across many current/external pairs.

\subsection{Balancing and reweighting methods}

A related class of methods approaches external-data use through balancing or
reweighting rather than hierarchical priors.
Propensity-score matching and weighting
\citep{rosenbaum1983central,stuart2010matching} adjust the external cohort so
that its covariate distribution better matches the current cohort.
Inverse-probability weighting is a standard implementation but can be
unstable when estimated propensity scores are close to 0 or 1
\citep{chesnaye2022introduction}, and overlap weighting
\citep{li2018balancing,li2019addressing} addresses this instability by
emphasizing subjects in the region of covariate overlap, with weights
proportional to $2\hat e(x)(1-\hat e(x))$.
Entropy balancing \citep{hainmueller2012entropy} reweights the external
cohort to match prespecified moments of the current cohort without estimating
a propensity model.
These ideas extend naturally to transportability and generalizability of
causal estimands across populations
\citep{dahabreh2019generalizing}.
Regulators have also begun to formalize the use of external and synthetic
controls in pivotal studies \citep{thorlund2020synthetic,fda2023external}.

A key limitation of weighting-based approaches, however, is that they
primarily address differences in baseline \emph{covariates} but not
differences in the \emph{outcome} generation mechanism
\citep{viele2014use,dahabreh2019generalizing}.
In Alzheimer's cohorts in particular, an external and a current source can
also differ in recruitment period, diagnostic practice, assessment schedule,
and baseline conversion risk.
In that case, balancing baseline covariates may not be enough: even a
well-calibrated external model may encode a shifted outcome relationship
relative to the current cohort.
This calibration-level mismatch, rather than covariate imbalance, is the
dominant failure mode in our ADNI case study.

\subsection{Neural posterior estimation and simulation-based inference}

Simulation-based inference (SBI) is a family of methods that perform
Bayesian inference using only the ability to simulate from the generative
model \citep{cranmer2020frontier}.
Neural posterior estimation (NPE) is the branch of SBI that trains a
conditional density estimator $q_{\phi}(\theta \mid x)$ directly on
simulated parameter--data pairs
\citep{papamakarios2016fast,papamakarios2019sequential,greenberg2019automatic,
lueckmann2021benchmarking}, with training objective
\begin{equation}
\mathcal{L}(\phi)
=
\mathbb{E}_{\theta \sim \pi(\theta),\, x \sim p(x\mid \theta)}
\left[-\log q_{\phi}(\theta \mid x)\right].
\end{equation}
Once trained, the network can produce an approximate posterior for a new
observation in a single forward pass, avoiding repeated posterior sampling.
Normalizing flows are the standard choice for $q_\phi$ because they provide
flexible, exactly-evaluable densities
\citep{papamakarios2021normalizing}, and open-source toolkits such as
\texttt{sbi} \citep{tejerocantero2020sbi} and \texttt{BayesFlow}
\citep{radev2022bayesflow} have made these methods practical for scientific
applications ranging from neuroscience \citep{goncalves2020training} to
epidemiology.

Two features of NPE are directly relevant to dynamic borrowing.
First, inference is \emph{amortized}: the upfront simulation-and-training
cost can be reused across many new dataset pairs, which matches the setting
where the same borrowing procedure must be evaluated repeatedly as cohorts
change.
Second, the network's input can be any summary of the two sources, not just
a single compatibility scalar---so long as the simulator covers the
relevant exchangeability regimes.
This expands the representational budget that governs how borrowing strength
reacts to different kinds of cohort shift.

For dynamic borrowing, the representation supplied to the posterior
estimator must therefore encode both information about the target estimand
and information about \emph{source compatibility}.
This requirement differs from many standard SBI settings, in which the
primary challenge is compressing data without losing information about a
single target parameter.
In borrowing problems, the representation must additionally preserve
structure that is informative about when external data should and should
not influence inference---the perspective we take in the remainder of the
paper.
